% Ley 1480 de 2011 - fragmento del corpus
\documentclass[11pt,a4paper]{article}
\usepackage[utf8]{inputenc}
\usepackage[T1]{fontenc}
\usepackage{lmodern}
\usepackage[spanish,es-noshorthands]{babel}
\usepackage[margin=2.5cm]{geometry}
\setlength{\parindent}{0pt}\setlength{\parskip}{0.7em}
\begin{document}
\section*{Ley 1480 de 2011}

\textit{FRAGMENTO ACADÉMICO — Avance 1}

Ley 1480 de 2011 — "Por medio de la cual se expide el Estatuto del Consumidor" (Colombia).

\textit{Nota: reproducción parcial y resumida con fines didácticos. Verifique siempre contra el texto oficial y vigente.}

\subsection*{Artículo 5. Definiciones.}

3. Consumidor: Toda persona natural o jurídica que, como destinatario final, adquiera, disfrute o utilice un determinado producto, cualquiera que sea su naturaleza, para la satisfacción de una necesidad propia, privada, familiar o doméstica y empresarial cuando no esté ligada intrínsecamente a su actividad económica.

5. Garantía: Obligación temporal, solidaria a cargo del productor y el proveedor, de responder por el buen estado del producto y la conformidad del mismo con las condiciones de idoneidad, calidad y seguridad legalmente exigibles o las ofrecidas.

\subsection*{Artículo 7. Garantía legal.}

Es la obligación, en los términos de esta ley, a cargo de todo productor y/o proveedor de responder por la calidad, idoneidad, seguridad y el buen estado y funcionamiento de los productos. En la prestación de servicios en el que el resultado es garantizado, la garantía estará dada por las condiciones de calidad e idoneidad ofrecidas al consumidor.

Ante los consumidores, la responsabilidad por la garantía legal recae solidariamente en los productores y proveedores respectivos.

\subsection*{Artículo 8. Términos de la garantía legal.}

El término de la garantía legal será el dispuesto por la ley o por la autoridad competente. A falta de disposición de obligatorio cumplimiento, será el anunciado por el productor y/o proveedor. El término de la garantía legal empezará a correr a partir de la entrega del producto al consumidor. De no indicarse el término, el mínimo será de un año para productos nuevos.

\subsection*{Artículo 11. Aspectos incluidos en la garantía legal.}

Corresponden a la garantía legal las siguientes obligaciones:

1. Presentado un defecto de calidad, idoneidad o seguridad de un producto, la reparación totalmente gratuita de los defectos y el suministro de los repuestos y mano de obra.

3. En los casos de reparación, se debe suministrar al consumidor un producto igual o de similares características, mientras se efectúa la reparación, cuando esta última tarde más de diez (10) días hábiles.

Cuando en un producto se hayan efectuado dos reparaciones por el mismo defecto y este se presente de nuevo, se procederá al cambio del bien por otro de las mismas características o a la devolución del dinero, a elección del consumidor.

\subsection*{Artículo 47. Retracto.}

En todos los contratos para la venta de bienes y prestación de servicios mediante sistemas de financiación otorgada por el productor o proveedor, venta de tiempos compartidos o ventas que utilizan métodos no tradicionales o a distancia, que por su naturaleza no deban consumirse o no hayan comenzado a ejecutarse antes de cinco (5) días, se entenderá pactado el derecho de retracto por parte del consumidor. En el evento en que se haga uso de la facultad de retracto, se resolverá el contrato y se deberá reintegrar el dinero que el consumidor hubiese pagado. El consumidor deberá devolver el producto al productor o proveedor por los mismos medios y en las mismas condiciones en que lo recibió. El término máximo para ejercer el derecho de retracto será de cinco (5) días hábiles contados a partir de la entrega del bien o de la celebración del contrato en caso de prestación de servicios. El productor o proveedor devolverá en dinero al consumidor todas las sumas pagadas, a más tardar dentro de los treinta (30) días calendario siguientes a aquel en que ejerció el derecho.

\subsection*{Artículo 50. Contenido de la información en comercio electrónico.}

Sin perjuicio de las demás obligaciones establecidas en la presente ley, cuando se trate de ventas mediante comercio electrónico, el proveedor deberá, entre otros: informar en todo momento de forma cierta, suficiente, clara, accesible y actualizada su identidad, las características del bien o servicio, el precio total incluidos los impuestos y costos de envío, el medio y la forma de pago, y el derecho de retracto que le asiste al consumidor y el procedimiento para ejercerlo.
\end{document}
